% Faithful native semantic reconstruction of physical page 15 / printed page 14.
\begin{HLAClauseSemanticPage}{14}
\HLAStableID{m3.page.015}
\HLAStableID{m3.p015.r01}
\HLAClauseRule{91.020}{693.000}{36.240}{0.480}
\HLAStableID{m3.p015.b01}
\HLAClausePlacedBlock{90.000}{69.646}{262.681}{%
  {\HLAArial\fontsize{9.960bp}{11.753bp}\selectfont\raggedright \resizebox{262.681bp}{!}{\mbox{\textbf{1.4.4 Relationship of HLA and object-oriented concepts}}}\par}
}
\HLAStableID{m3.p015.b02}
\HLAClausePlacedBlock{89.999}{93.785}{432.033}{%
  {\HLATimes\fontsize{9.960bp}{12.000bp}\selectfont\raggedright \resizebox{431.944bp}{!}{\mbox{Although the HLA OMT is the standardized documentation structure for HLA object models, FOMs, FOM}}\newline{}\resizebox{431.952bp}{!}{\mbox{modules, SOMs, and SOM modules do not completely correspond to common definitions of object models}}\newline{}\resizebox{431.984bp}{!}{\mbox{in object-oriented analysis and design (OOAD) techniques. In the OOAD literature, an object model is}}\newline{}\resizebox{431.997bp}{!}{\mbox{described as an abstraction of a system developed for the purpose of fully understanding the system. To}}\newline{}\resizebox{431.929bp}{!}{\mbox{achieve this understanding, most object-oriented (OO) techniques recommend defining several views of the}}\newline{}\resizebox{431.986bp}{!}{\mbox{system. For HLA object models, the intended scope of the system description is much narrower, focusing}}\newline{}\resizebox{431.960bp}{!}{\mbox{mainly on requirements and capabilities for federate information exchange. For SOMs and SOM modules,}}\newline{}\resizebox{431.980bp}{!}{\mbox{the intent is to describe the public interface of the federate in terms of an identified set of supported HLA}}\newline{}\resizebox{431.975bp}{!}{\mbox{object classes and interaction classes. A more complete description of how a federate is designed and}}\newline{}\resizebox{431.987bp}{!}{\mbox{functions internally (for example, a traditional OO object model) should be provided via documentation}}\newline{}\resizebox{432.033bp}{!}{\mbox{resources other than the SOM. For FOMs and FOM modules, the intent is to describe information exchange}}\newline{}\resizebox{175.081bp}{!}{\mbox{that happens during a federation execution.}}\par}
}
\HLAStableID{m3.p015.b03}
\HLAClausePlacedBlock{89.999}{249.783}{432.080}{%
  {\HLATimes\fontsize{9.960bp}{12.000bp}\selectfont\raggedright \resizebox{432.080bp}{!}{\mbox{Differences between HLA and OOAD principles and concepts also appear at the individual object level. In}}\newline{}\resizebox{431.978bp}{!}{\mbox{the OOAD literature, objects are defined as software encapsulations of data and operations (methods). In the}}\newline{}\resizebox{431.939bp}{!}{\mbox{HLA, objects are defined entirely by the identifying characteristics (attributes), values of which are}}\newline{}\resizebox{432.017bp}{!}{\mbox{exchanged between federates during a federation execution. Any OO-related behaviors and operations that}}\newline{}\resizebox{432.024bp}{!}{\mbox{affect the values of HLA object attributes are kept resident in the federates. Although HLA class structures}}\newline{}\resizebox{431.991bp}{!}{\mbox{are driven by subscription requirements and FOM growth concerns, class structures in OO systems are}}\newline{}\resizebox{240.972bp}{!}{\mbox{typically driven by efficiency and maintainability concerns.}}\par}
}
\HLAStableID{m3.p015.b04}
\HLAClausePlacedBlock{89.998}{345.781}{432.105}{%
  {\HLATimes\fontsize{9.960bp}{12.000bp}\selectfont\raggedright \resizebox{432.103bp}{!}{\mbox{As discussed above, HLA object classes are described by the attributes that are defined for them. These are,}}\newline{}\resizebox{432.045bp}{!}{\mbox{in OO parlance, data members of the class. These attributes, which are abstract properties of HLA object}}\newline{}\resizebox{432.032bp}{!}{\mbox{classes, are referred to as class attributes. HLA object instances are spawned via an HLA service using an}}\newline{}\resizebox{432.003bp}{!}{\mbox{HLA object class as a template. Each attribute contained by an HLA object instance is called an\space{}\textit{instance}}}\newline{}\resizebox{36.930bp}{!}{\mbox{\textit{attribute}.}}\par}
}
\HLAStableID{m3.p015.b05}
\HLAClausePlacedBlock{89.998}{417.780}{432.124}{%
  {\HLATimes\fontsize{9.960bp}{12.000bp}\selectfont\raggedright \resizebox{432.030bp}{!}{\mbox{OO objects interact via message passing, in which one OO object invokes an operation provided by another}}\newline{}\resizebox{432.124bp}{!}{\mbox{OO object. HLA objects do not directly interact. It is the federates that interact, via HLA services, by}}\newline{}\resizebox{431.972bp}{!}{\mbox{updating instance attribute values or sending interactions. Also, responsibility for updating the instance}}\newline{}\resizebox{431.906bp}{!}{\mbox{attributes associated with an HLA object instance can be distributed among different federates in a}}\newline{}\resizebox{431.976bp}{!}{\mbox{federation (effectively distributing responsibility for maintaining the HLA object instance’s state across the}}\newline{}\resizebox{432.013bp}{!}{\mbox{federation), whereas OO objects encapsulate state locally and associate update responsibilities with}}\newline{}\resizebox{386.098bp}{!}{\mbox{operations that are closely tied to the object’s implementation in an OO programming language.}}\par}
}
\HLAStableID{m3.p015.b06}
\HLAClausePlacedBlock{89.998}{513.779}{432.032}{%
  {\HLATimes\fontsize{9.960bp}{12.000bp}\selectfont\raggedright \resizebox{431.974bp}{!}{\mbox{Specialization of HLA object and interaction classes are discussed in terms of data tree structures. A data}}\newline{}\resizebox{431.988bp}{!}{\mbox{tree structure is defined as, “an abstract hierarchical structure consisting of nodes [classes] connected by}}\newline{}\resizebox{432.028bp}{!}{\mbox{branches, in which; each branch connects one node [class] to a directly subsidiary node; there is a unique}}\newline{}\resizebox{431.960bp}{!}{\mbox{node called the root [root class] which is not subsidiary to any other node; and every node [class] besides the}}\newline{}\resizebox{432.032bp}{!}{\mbox{root is directly subsidiary to exactly one other node [class]” (\textit{The IEEE Standards Dictionary Online}).\HLATextSup{8}}}\par}
}
\HLAStableID{m3.p015.b07}
\HLAClausePlacedBlock{90.000}{573.786}{329.616}{%
  {\HLATimes\fontsize{9.960bp}{11.753bp}\selectfont\raggedright \resizebox{329.616bp}{!}{\mbox{Furthermore, a node without any subsidiary node is called a leaf node [leaf class].}}\par}
}
\HLAStableID{m3.p015.b08}
\HLAClausePlacedBlock{89.999}{597.785}{432.046}{%
  {\HLATimes\fontsize{9.960bp}{12.000bp}\selectfont\raggedright \resizebox{432.045bp}{!}{\mbox{In addition to the stated semantic variations in shared terminology, other differences may also exist. Precise}}\newline{}\resizebox{217.822bp}{!}{\mbox{definitions of all HLA terms can be found in Clause 3.}}\par}
}
\HLAStableID{m3.p015.b09}
\HLAClausePlacedBlock{90.000}{696.453}{432.016}{%
  {\HLATimes\fontsize{7.980bp}{9.000bp}\selectfont\raggedright \resizebox{432.016bp}{!}{\mbox{\HLATextSup{8}\textit{IEEE Standards Dictionary Online}\space{}is available at\space{}\textcolor{HLAc0000FF}{\uline{http://dictionary.ieee.org/}}. An IEEE account is required for access to the dictionary,}}\newline{}\resizebox{215.384bp}{!}{\mbox{and one can be created at no charge on the dictionary sign-in page.}}\par}
}
\end{HLAClauseSemanticPage}
